% Appendix, from exam/README.md: what the two written papers are for, what they cover, the
% conventions they assume, and how they are marked.

\chapter{The Self-Assessment Papers}
\label{exam:ch:about}

This appendix and the four after it hold two complete three-hour papers for \emph{Embedded C},
with worked solutions. Appendices~\ref{exa:ch:paper} and \ref{exb:ch:paper} are the papers,
questions only; Appendices~\ref{ansa:ch:answers} and \ref{ansb:ch:answers} are their model answers,
with marks.

Each paper is self-contained. It states its own rubric, and it closes with a register reference
covering every ATmega328P register and bit it refers to, so nobody sitting one needs a datasheet,
a chapter of this book, or this appendix.

\section*{What these are for}
The book has \textbf{no exam and nothing is marked}. What you get instead is the exercise set after
every chapter, and the questions in every chapter's review.

These papers do not replace that. \textbf{They are here for you to test your own skills and
knowledge, nothing more.} They gate nothing, they are not a qualification, and no part of the book
requires them. Nothing in the companion repository depends on them, and neither \code{make build}
nor \code{make format-check} knows they exist.

What they check is something the exercises cannot. Every exercise is done at a keyboard, with the
chapter open beside you, the datasheet beside that, and a board on the desk to tell you whether it
worked. These papers give you a page of register names and take the rest away, then ask what you
can reconstruct --- which of the rules you have absorbed and which you have merely been looking up.

\textbf{Take one after the last chapter.} Both papers draw on all ten chapters, so sitting one
partway through examines material you have not read yet, and the result says more about how far
you have read than about what you have understood.

\section*{The two papers}
Both cover the whole book, Chapters~1 to 10, at the same weighting. They share no question. Either
can be used alone; use both as a main sitting and a resit, or in alternate years.

\begin{booktable}{@{}p{4.6em}Lp{5.4em}r@{}}
  \thead{Question} & \thead{Topic} & \thead{Chapters} & \thead{Marks} \\
  \midrule
  1 & The C language underneath: pointers, macros, types & 1 & 12 \\
  2 & Registers and bitwise I/O & 2 & 12 \\
  3 & Interrupts, edges and \code{volatile} & 3 & 12 \\
  4 & The ADC, and timers that count & 4, 5 & 13 \\
  5 & Serial communication & 6 & 13 \\
  6 & EEPROM and the watchdog & 7, 8 & 13 \\
  7 & Structs as drivers & 9 & 12 \\
  8 & Encapsulation and polymorphism in C & 10 & 13 \\
  \midrule
    & & & \textbf{100} \\
\end{booktable}

\textbf{Both papers mix theory with code}, because the subject does. Roughly half the marks are for
prose --- what \code{volatile} guarantees, why an ISR must not transmit, what an opaque struct
protects --- and roughly half are for reading or writing C on paper.

The two papers differ in which half they lean on:
\begin{itemize}
  \item \textbf{Paper A leans towards code that is put in front of you.} All eight of its questions
    hand you a snippet that compiles and is wrong: a macro whose missing parentheses silently turn
    50 into 32, a \code{free()} that leaves the caller holding the freed address, a \code{setup()}
    that writes board pin numbers where port bit positions belong and therefore writes zero to
    every register, an interrupt armed at the peripheral but never enabled globally, an ADC read
    with no reference and no wait, a UART driver with the baud \emph{rate} written into
    \code{UBRR0}, an EEPROM write with \code{EEPE} set before \code{EEMPE}, a watchdog ``pet'' that
    sets a bit instead of executing an instruction, a \code{gpio_toggle()} that uses \code{^=} on
    \code{PINx}, and a vtable whose interface field is not the struct's first member. Naming the
    defect is half the marks; saying what it does to the running program is the other half.
  \item \textbf{Paper B leans towards code you write and design decisions you justify.} It asks for
    a complete three-LED chase program written from nothing, an enumeration and its validity check,
    a \code{realloc()}-based resize that survives a failed allocation, an INT1 handler with an
    atomic toggle, \code{adc_read()} in full, \file{driver/serial.h} in full with its Doxygen, the
    EEPROM read/write pair, the watchdog timeout enumeration and its two odd values, a GPIO driver
    that maps an Arduino pin number to a port, a pin reservation table, and a \code{gpio_stub_t}
    that satisfies the vtable interface with no hardware behind it. Several parts are pure
    judgement: which of four values belongs in EEPROM and why, where a report belongs relative to
    the interrupt that triggered it, and what opacity takes away from a caller.
\end{itemize}

Neither paper asks anything of a compiler that a careful reader could not answer from the
chapters. Where a question turns on a rule that is easy to state and easy to get backwards ---
which register reports a pin's actual level, which of \code{EEMPE} and \code{EEPE} comes first, why
\code{PINx = (1U << n)} is atomic and \code{PORTx ^= (1U << n)} is not --- the solution says so at
length, because that is where the marks and the understanding both are.

\section*{Conventions the papers assume}
Both papers state these in their own rubric, so a candidate never has to have read this appendix.
\begin{itemize}
  \item \textbf{C11 on an ATmega328P at 16\,MHz}, built with \code{avr-gcc}. Standard and AVR
    headers may be assumed included unless a question asks for them.
  \item \textbf{Register and bit names are not memory-tested.} Each paper closes with a register
    reference of its own, so nobody sitting one needs a datasheet in the room. Misspelling a bit
    costs nothing; reaching for the wrong \emph{register}, or the right register at the wrong
    \emph{time}, costs everything. What the reference deliberately withholds is the ordering rules,
    the timing rules and the formulas --- those are the parts being examined.
  \item \textbf{Code is marked on semantics, not syntax.} A missing semicolon, a forgotten
    \code{#include} or a brace in the wrong column costs nothing. A missing \code{volatile} costs
    everything.
  \item \textbf{The book's house style is expected} where a question asks for it:
    \code{snake_case} with \code{UPPER_CASE} macros, \code{static} on anything private, header
    guards named after the file, Doxygen in the imperative mood, \code{self} as a pointer (and
    \code{const} when the function only reads), Yoda comparisons, braces on every \code{if},
    fixed-width types with \code{U} suffixes, \code{const} on by-value parameters in the definition
    but not the prototype, and \code{volatile} on anything shared with an ISR.
  \item \textbf{``State the consequence'' means the consequence.} Naming a defect earns half the
    marks; saying what it does to the running program earns the other half. ``This is wrong''
    scores nothing.
  \item \textbf{Where a question asks for a stated number of defects}, listing more is not
    penalised, but only that many are marked, so the strongest answers go first.
\end{itemize}

\section*{Marking}
Every solution is written to be marked by somebody who has read the chapters and does not otherwise
write embedded C for a living, so each carries the reasoning rather than the answer alone. Marks
are shown per part, and the arithmetic in the timer, ADC and baud-rate questions is worked through
step by step rather than asserted.

Three conventions worth agreeing before a paper is marked:
\begin{itemize}
  \item \textbf{Method carries the marks.} A correct diagnosis with a clumsy fix is worth more than
    a correct fix with no diagnosis. Several questions feed the next part, so an error should be
    followed through rather than penalised twice.
  \item \textbf{The named traps are worth full marks on their own.} A number of parts exist entirely
    to see whether a candidate walks into a specific mistake: reading a button from \code{PORTx}, a
    \code{1U << 9U} written to an 8-bit register, a plain \code{=} where \code{|=} was needed, an
    ISR-shared flag without \code{volatile}, \code{EEPE} before \code{EEMPE}, a watchdog pet written
    as a register access, \code{^=} on \code{PINx}, a \code{gpio_t} passed by value, an interface
    field that is not first in its struct. Where a solution flags one of these, a candidate who
    walks into it loses those marks and no others.
  \item \textbf{The discussion parts are not decoration.} ``State why'' is where the book's actual
    content is, and a paper marked only on whether the code would compile would pass a candidate
    who has understood none of it.
\end{itemize}
